\documentclass[12pt,a4paper]{article}
\usepackage[margin=25mm, headheight=15pt, headsep=10pt]{geometry}
\usepackage{fontspec}
\usepackage[table]{xcolor} % Note: table option is required for rowcolors
\usepackage{titlesec}
\usepackage{tocloft}
\usepackage{fancyhdr}
\usepackage{booktabs}
\usepackage{tcolorbox}
\tcbuselibrary{skins, breakable}
\usepackage[colorlinks=true, linkcolor=emerald, urlcolor=emerald, bookmarks=true]{hyperref}

\definecolor{emerald}{RGB}{0,102,51}
\definecolor{darkred}{RGB}{139,0,0}
\definecolor{lightgray}{RGB}{245,245,245}

\usepackage[nil]{babel}
\babelprovide[import, main]{bengali}
\babelprovide[import]{arabic}
\babelfont{rm}[Renderer=HarfBuzz,Scale=1.0]{Noto Serif Bengali}
\babelfont{sf}[Renderer=HarfBuzz]{Noto Sans Bengali}
\babelfont{tt}[Renderer=HarfBuzz]{Noto Sans Bengali}
\babelfont[arabic]{rm}[Renderer=HarfBuzz,Scale=1.1]{Noto Naskh Arabic}

% Section styling
\titleformat{\section}{\Large\bfseries\color{emerald}}{\thesection.}{0.5em}{}
\titleformat{\subsection}{\large\bfseries\color{emerald!85!black}}{\thesubsection.}{0.5em}{}
\titleformat{\subsubsection}{\normalsize\bfseries\color{emerald!70!black}}{\thesubsubsection.}{0.5em}{}
\titlespacing*{\section}{0pt}{18pt}{8pt}

% Fancy Header/Footer setup
\pagestyle{fancy}
\fancyhf{} % clear all header and footer fields
\fancyhead[L]{\color{emerald}\small\textbf{\nouppercase{\leftmark}}}
\fancyfoot[L]{\scriptsize\color{black!50}{\lat{AI}}-সংকলিত, আলেম-যাচাইকৃত নয়}
\fancyfoot[C]{\color{emerald!70!black}\thepage}
\renewcommand{\headrulewidth}{0.4pt}
\renewcommand{\headrule}{\hbox to\headwidth{\color{emerald}\leaders\hrule height \headrulewidth\hfill}}
\renewcommand{\footrulewidth}{0pt}

% Custom boxes
% 1. Ayat/Hadith Box
\newtcolorbox{ayatbox}[1][]{
    enhanced, breakable, 
    colback=emerald!5, 
    colframe=emerald, 
    boxrule=0.5pt, 
    arc=3pt, 
    left=10pt, right=10pt, top=10pt, bottom=10pt,
    #1
}

% 2. Quote/Translation Box
\newtcolorbox{quotebox}[1][]{
    enhanced, breakable,
    frame hidden,
    colback=lightgray,
    borderline west={3pt}{0pt}{emerald},
    left=15pt, right=10pt, top=10pt, bottom=10pt,
    sharp corners,
    #1
}

% 3. AI-generated notice box
\newtcolorbox{warnbox}[1][]{
    enhanced, breakable,
    colback=red!4,
    colframe=darkred,
    boxrule=0.8pt,
    arc=3pt,
    left=10pt, right=10pt, top=10pt, bottom=10pt,
    #1
}

% Commands
\newcommand{\arab}[1]{\foreignlanguage{arabic}{#1}}
\newcommand{\kib}[1]{\textcolor{emerald}{\textbf{#1}}}

% Latin (English) fallback: Noto Serif Bengali has NO Latin glyphs
\newfontfamily\latfont[Renderer=HarfBuzz]{Noto Serif}
\newcommand{\lat}[1]{{\latfont #1}}

\setlength{\parskip}{5pt}
\setlength{\parindent}{0pt}

% Fix for missing bullet in Noto Serif Bengali
\renewcommand{\labelitemi}{$\bullet$}



\begin{document}
\begin{center}{\Large\bfseries\color{emerald} ব্যক্তিগত উন্নতির জন্য কাজের বিষয় — মুহাসাবা, দোয়া ও অনুশীলন}\end{center}
\vspace{1em}
\section*{ব্যক্তিগত উন্নতির জন্য কাজের বিষয় — মুহাসাবা, দোয়া ও অনুশীলন}

\begin{warnbox}
 \textbf{গুরুত্বপূর্ণ নোটিশ (\lat{Important} \lat{Notice}):} এই কন্টেন্টটি \lat{AI} (কৃত্রিম বুদ্ধিমত্তা) দিয়ে প্রস্তুত ও সংকলিত। \lat{AI} কঠোর সূত্র-মান নীতি মেনে শুধু নির্ভরযোগ্য উৎস থেকে তথ্য সংগ্রহ করেছে — কেবল সহীহ/হাসান হাদিস ও সহীহ সনদের আসার রাখা হয়েছে, দুর্বল সনদ বাদ দেওয়া হয়েছে। তবে এটি কোনো আলেম বা মুহাদ্দিস কর্তৃক যাচাইকৃত নয়।
\end{warnbox}


\begin{quote}
এই ফাইলটি কিবর ও উজব থেকে বাঁচার \textbf{\lat{practical} (ব্যবহারিক)} পথ দেখায় — প্রতিদিনের মুহাসাবা (আত্মনিরীক্ষা), \lat{dua} (দোয়া) এবং সাহাবিদের \lat{character} (চরিত্র) অনুশীলনের \lat{technique} (কৌশল)।
\end{quote}

\medskip\hrule\medskip

\subsection*{১. প্রতিদিনের মুহাসাবা (আত্মনিরীক্ষা)}

\begin{quote}
উমার (রাঃ) বলতেন: \textbf{\lat{“}তোমাদের বিচার হওয়ার আগে নিজেরা নিজেদের বিচার করো। তোমাদের আমল ওজন হওয়ার আগে নিজেরা ওজন করো। আর বড় হিসাবের (কিয়ামত) জন্য প্রস্তুত থাকো।\lat{”}}
\end{quote}

\textbf{প্রতিদিন সকাল-সন্ধ্যা ৫ মিনিট নিজেকে এই প্রশ্নগুলো করুন — এটি \lat{daily} \lat{self-check} (প্রতিদিনের আত্মপরীক্ষা):}

\begin{enumerate}
\item আজ আমার মনে কোথায় কিবর বা উজব জেগেছিল — কোনো \lat{pride} (গর্ব) বা \lat{self-satisfaction} (আত্মতুষ্টি)?
\item কেউ আমাকে প্রশংসা করলে আমার মনের অবস্থা কী ছিল — \lat{grateful} (কৃতজ্ঞ) নাকি \lat{conceited} (দাম্ভিক)?
\item আমি কি কাউকে হেয় দৃষ্টিতে দেখেছি বা তুচ্ছ করেছি — কোনো \lat{look-down} (তুচ্ছ দৃষ্টিভঙ্গি)?
\item আমার আমল (নামাজ, রোজা, জাকাত, সদকা) কি একমাত্র আল্লাহর জন্য ছিল — \lat{pure} \lat{sincerity} (বিশুদ্ধ ইখলাস)?
\item কেউ আমাকে সৎপরামর্শ দিলে আমি তা গ্রহণ করেছি কি — না \lat{reject} (প্রত্যাখ্যান) করেছি?
\end{enumerate}
\medskip\hrule\medskip

\subsection*{২. কিবর ও উজব দূর করার দোয়া}

\textbf{সাধারণ দোয়া:}
\begin{quote}
\textbf{\lat{“}হে আল্লাহ! আমাকে কিবর, উজব, রিয়া, হাসাদ ও অন্যান্য খারাপ চরিত্র থেকে রক্ষা করুন। আমাকে বিনীত, নম্র ও সত্যিকারের মুমিন বানান। আমিন।\lat{”}}
\end{quote}

\textbf{নবীজির শিক্ষানো দোয়া (সহীহ):}
\begin{quote}
\textbf{\lat{“}হে আল্লাহ! আমাকে সুন্দর চরিত্রের পথ দেখান, আপনি ছাড়া আর কেউ তা দেখাতে পারে না। আমাকে মন্দ চরিত্র থেকে দূরে রাখুন, আপনি ছাড়া আর কেউ তা দূরে রাখতে পারে না।\lat{”}} — সুনানে নাসাই ৮৬১ (সহীহ)
\end{quote}

\medskip\hrule\medskip

\subsection*{৩. জ্ঞান প্রার্থনায় ইখলাস}

\begin{quote}
\textbf{শেখ ইবনে উসাইমীন (রহ.):} \textit{\lat{“}জ্ঞান একমাত্র আল্লাহর জন্য, মানুষের প্রশংসার জন্য নয়। জ্ঞান দিয়ে আল্লাহর সন্তুষ্টি চাও, মানুষের মন জয় নয়।\lat{”}}
\end{quote}

\begin{itemize}
\item পড়াশোনার শুরুতে \lat{intention} (নিয়ত) করুন: \textbf{এই জ্ঞান (\lat{knowledge}) আল্লাহর জন্য, যাতে আমি তার বিধান বুঝি ও আমল (\lat{practice}) করি।}
\item যে জ্ঞান আলেমদের সামনে দেখানোর, \lat{debate} (বিতর্ক) জেতার বা চাকরির সিঁড়ি বানানোর হয় — সে জ্ঞান উজব-রিয়ার \lat{seed} (বীজ)।
\item নিয়মিত নিয়তের পর্যালোচনা (মুরাকাবা) করুন — নিয়ত নষ্ট হলে নতুন করে \lat{fresh} \lat{intention} (নতুন নিয়ত) নিন।
\end{itemize}
\medskip\hrule\medskip

\subsection*{৪. সাহাবিদের পথে চলা — অনুশীলনের সারণি}

\begin{table}[h]\centering\small
\begin{tabular}{lll}
গুণ & সাহাবীর উদাহরণ & অনুশীলন \\
\hline
বিনয় & নবীজি, আবু বকর, উমার & দরিদ্র/দাসের সাথে মাটিতে বসে খাওয়া, নিজের কাজ নিজে করা — \lat{humility} \lat{practice} (বিনয়ের অনুশীলন) \\
মুহাসাবা & উমার, হানজালা & প্রতিদিন আত্মনিরীক্ষা, গুনাহ মনে হলে কান্না — \lat{self-accountability} (আত্মজবাবদিহি) \\
সাদগী & উমার, আবু বকর & সাধারণ পোশাক, বিলাসিতা পরিহার — \lat{simplicity} (সরলতা) \\
সত্যের প্রতি নমনীয়তা & আবু যার, ইবনে মাসউদ & ভুল হলে মেনে নেওয়া, লজ্জিত হয়ে সংশোধন — \lat{openness} \lat{to} \lat{truth} (সত্যের প্রতি উন্মুক্ততা) \\
\hline
\end{tabular}
\end{table}

\medskip\hrule\medskip

\subsection*{৫. সত্যিকারের প্রতিকার — মনে রাখার মূলনীতি}

\begin{quote}
\textbf{\lat{“}কিবর হলো সত্যকে তুচ্ছ করা ও মানুষকে হেয় করা।\lat{”}} — এই সংজ্ঞাটি মুহূর্তে মুহূর্তে নিজের মনে আনুন।
\end{quote}

\begin{enumerate}
\item \textbf{সত্যকে তুচ্ছ না করা:} কেউ \lat{truth} (সত্য) বললে — ছোট, দরিদ্র বা কম জ্ঞানী হলেও — তা মেনে নিন।
\item \textbf{মানুষকে হেয় না করা:} কেউ \lat{mistake} (ভুল) করলে তাকে ছোট না করে পরামর্শ দিন — নম্রভাবে \lat{correct} (সংশোধন) করুন।
\item \textbf{আত্মতুষ্টি না করা:} নিজের কোনো গুণ বা আমলের জন্য নিজেকে ধন্য মনে করবেন না — সবই আল্লাহর \lat{blessing} (নেয়ামত)।
\item \textbf{প্রশংসা শুনলে:} মনে রাখুন প্রশংসা একটি \lat{test} (পরীক্ষা); সাথে সাথে আল্লাহর কাছে \lat{humility} (বিনয়)-এর দোয়া করুন।
\item \textbf{নিন্দা শুনলে:} রাগ (\lat{anger}) নয়, \lat{self-reflection} (আত্মবিশ্লেষণ) করুন — নিন্দায় সত্য থাকলে সংশোধন হোন।
\end{enumerate}
\medskip\hrule\medskip

\subsection*{৬. সমাপনী দোয়া}

\begin{quote}
\textbf{\lat{“}হে আল্লাহ! আমাদের অন্তর থেকে কিবর ও উজব দূর করুন। আমাদের বিনীত, নম্র, সত্যিকারের মুমিন এবং আপনার প্রিয় বান্দা বানান। আমিন, ইয়া রাব্বাল আলামীন।\lat{”}}
\end{quote}


\end{document}
